% ==================================================================
\section{Iterative: Tile Task Decomposition}
\label{sec:iter-tile}
% ==================================================================

\subsection{Implementation}
\label{sec:iter-tile-impl}

The \textbf{original tile iterative} strategy maps the $\mathit{ROWS}\times\mathit{COLS}$ image
onto a grid of square tiles of size $\mathit{TILE}\times\mathit{TILE}$ pixels.
The outer double loop iterates over tile origins with stride \texttt{TILE}; each iteration
body is wrapped in a single \texttt{\#pragma~omp~task}, so all tiles are eligible for
parallel execution.
The task takes \texttt{x} and \texttt{y} as \texttt{firstprivate} to avoid the data-race
that would arise if the loop variables were shared across concurrently executing tasks.

The parallel region is opened with \texttt{\#pragma~omp~parallel} followed immediately by
\texttt{\#pragma~omp~single}, so only one thread generates the tasks while all others
are available to execute them.

\begin{lstlisting}[caption={Key parallel section of \texttt{mandel-omp-iter.c} (Tile strategy).},
                   label={lst:iter-tile}]
#pragma omp parallel
#pragma omp single
for (int y = 0; y < ROWS; y += TILE)
    for (int x = 0; x < COLS; x += TILE) {
        #pragma omp task firstprivate(x, y)
        {
            int equal = 1;
            // --- Check vertical TILE borders ---
            int ref = pixel_dwell(..., x, y, ...);
            for (int py = y; py < y + TILE; py++) {
                M[py][x]          = pixel_dwell(..., x,          py, ...);
                M[py][x+TILE-1]   = pixel_dwell(..., x+TILE-1,  py, ...);
                equal &= (ref == M[py][x]);
                equal &= (ref == M[py][x+TILE-1]);
            }
            // --- Check horizontal TILE borders ---
            for (int px = x; px < x + TILE; px++) {
                M[y][px]          = pixel_dwell(..., px, y,         ...);
                M[y+TILE-1][px]   = pixel_dwell(..., px, y+TILE-1, ...);
                equal &= (ref == M[y][px]);
                equal &= (ref == M[y+TILE-1][px]);
            }
            if (equal) {
                fill_tile(M, y, x, TILE, M[y][x]);   // uniform colour
            } else {
                compute_tile(M, y, x, TILE, ...);     // full computation
            }
        }
    }
\end{lstlisting}

Dependencies and data-sharing between tasks are avoided by construction: each task writes
only to the memory region $[y,\,y+\mathit{TILE})\times[x,\,x+\mathit{TILE})$, which is
disjoint from every other task's region.
Histogram updates use \texttt{\#pragma~omp~atomic} to protect the shared counter array.

\subsection{Modelfactor Analysis}
\label{sec:iter-tile-mf}

% TODO: Run  sbatch submit-strong-extrae.sh mandel-omp-iter
%       and paste the Modelfactors tables here.

\todo[inline]{Insert Modelfactors table for \texttt{mandel-omp-iter} (16 threads) and
identify the main bottlenecks (e.g.\ scheduling overhead, load imbalance, synchronisation).}

% Example skeleton (replace values with real data):
% \begin{table}[H]
% \centering
% \caption{Modelfactors for Iterative Tile with 16 threads.}
% \begin{tabular}{lr}
% \toprule
% Factor & Value \\
% \midrule
% Useful computation & \% \\
% Task scheduling overhead & \% \\
% Synchronisation & \% \\
% \bottomrule
% \end{tabular}
% \end{table}
%
% Main issues observed:
% - ...

\subsection{Paraver Analysis}
\label{sec:iter-tile-paraver}

% TODO: Open the Paraver trace generated for 16 threads and include:
%   1. The "Useful/Instantaneous parallelism" view.
%   2. The "OpenMP tasking / explicit tasks function created and executed" view.

\todo[inline]{Insert Paraver screenshots (16 threads) and comment on the task creation
pattern, load distribution, and any visible idle periods.}

% \begin{figure}[H]
%   \centering
%   \includegraphics[width=\linewidth]{fotos/iter-tile-paraver-parallelism.png}
%   \caption{Useful/Instantaneous parallelism -- Iterative Tile, 16 threads.}
%   \label{fig:iter-tile-paraver}
% \end{figure}

\subsection{Strong Scalability}
\label{sec:iter-tile-scalability}

% TODO: Run  sbatch submit-strong-omp.sh  and include the speedup plot.

\todo[inline]{Insert the strong scalability (speedup vs.\ number of threads) plot for
\texttt{mandel-omp-iter} and explain the observed trend.}

% \begin{figure}[H]
%   \centering
%   \includegraphics[width=0.75\linewidth]{fotos/iter-tile-speedup.png}
%   \caption{Strong scalability of Iterative Tile.}
%   \label{fig:iter-tile-speedup}
% \end{figure}
